La arquitectura híbrida DEB+GP requiere un mecanismo formal para cuantificar la discrepancia entre la predicción mecanicista y la realidad sensorial. Dado que el modelo DEB estima una tasa de producción de \ce{CO2} (masa por unidad de tiempo), mientras que el sensor NDIR registra una concentración acumulada en el headspace del reactor, no es posible comparar directamente ambas señales en el dominio instantáneo. En su lugar, se define una función de pérdida que opera sobre las \emph{pendientes de emisión} en ventanas temporales cerradas, donde la integración numérica del modelo DEB y la derivada numérica de la serie sensorial convergen a una misma magnitud física: tasa de cambio de concentración.

\subsection{Comparación de pendientes de emisión en ventanas cerradas}
\label{subsec:cap4_pendientes_ventanas}

Sea $C(t)$ la serie temporal de concentración de \ce{CO2} medida por el sensor NDIR, muestreada a \num{1} min de resolución. Para una ventana temporal cerrada de ancho $\Delta t_w = \num{30}$ min, centrada en el instante $t_k$, se define la pendiente observada mediante regresión lineal ordinaria:
\begin{equation}
    \label{eq:cap4_pendiente_sensor}
    \left.\frac{\Delta C}{\Delta t}\right|_{t_k} =
    \frac{\sum_{j=1}^{n_w}\bigl(t_j - \bar{t}\bigr)\bigl(C_j - \bar{C}\bigr)}
    {\sum_{j=1}^{n_w}\bigl(t_j - \bar{t}\bigr)^{2}},
\end{equation}
donde $n_w = \num{30}$ son los puntos dentro de la ventana, $\bar{t}$ y $\bar{C}$ son las medias temporales y de concentración, respectivamente, y $t_j \in [t_k - \Delta t_w/2,\; t_k + \Delta t_w/2]$. La elección de \num{30} min como ancho de ventana balancea dos requerimientos contrapuestos: por un lado, ventanas demasiado cortas amplifican el ruido de medición y la variabilidad de origen turbulenta en el headspace; por otro, ventanas excesivamente largas suavizan eventos transitorios de interés metabólico, como picos post-alimentación o transiciones térmicas \parencite{rossiEstimatingDynamicsGreenhouse2024}.

Simultáneamente, el modelo DEB —alimentado por el pseudo-observable $\hat{B}(t)$ del GP— produce la tasa teórica de emisión $\dot{m}_{\mathrm{CO_2}}^{\mathrm{DEB}}(t)$ en g/min. Para obtener una magnitud comparable con la pendiente sensorial, es necesario convertir la tasa másica en tasa de cambio de concentración. Asumiendo un reactor de volumen de cabeza $V_h$ (en L) conocido y una temperatura de operación $T$ tal que la densidad molar del gas se aproxima por el gas ideal, se tiene:
\begin{equation}
    \label{eq:cap4_conversion_masa_ppm}
    \dot{C}_{\mathrm{DEB}}(t) =
    \frac{\dot{m}_{\mathrm{CO_2}}^{\mathrm{DEB}}(t)}{V_h}
    \cdot \frac{R\,T}{M_{\mathrm{CO_2}}}
    \cdot 10^{6},
\end{equation}
donde $R = \num{8.314}$ J/(mol·K) es la constante universal de los gases, $M_{\mathrm{CO_2}} = \num{44.01}$ g/mol la masa molar del \ce{CO2}, y el factor $10^{6}$ convierte la fracción molar a ppm. En la práctica experimental de esta tesis, el reactor opera a $T \approx \num{303}$ K (\num{30} °C), de modo que el factor de conversión es aproximadamente constante durante el ciclo.

La media de la tasa modelo sobre la misma ventana de \num{30} min se calcula mediante integración numérica:
\begin{equation}
    \label{eq:cap4_media_modelo_ventana}
    \bar{\dot{C}}_{\mathrm{DEB}}(t_k) =
    \frac{1}{\Delta t_w}\int_{t_k - \Delta t_w/2}^{t_k + \Delta t_w/2}
    \dot{C}_{\mathrm{DEB}}(\tau)\,\mathrm{d}\tau
    \approx \frac{1}{n_w}\sum_{j=1}^{n_w}\dot{C}_{\mathrm{DEB}}(t_j).
\end{equation}

Con ambas magnitudes expresadas en ppm/min, se define el residuo operativo en la ventana $k$:
\begin{equation}
    \label{eq:cap4_residuo_operativo}
    \delta(t_k) =
    \left.\frac{\Delta C}{\Delta t}\right|_{t_k} -
    \bar{\dot{C}}_{\mathrm{DEB}}(t_k).
\end{equation}

La función de pérdida global del sistema, sobre un conjunto de $K$ ventanas no solapadas que cubren el ciclo experimental, es la suma de residuos cuadráticos ponderados por la incertidumbre del GP:
\begin{equation}
    \label{eq:cap4_funcion_perdida}
    \mathcal{L}(\boldsymbol{\theta}_{\mathrm{libre}}) =
    \sum_{k=1}^{K}
    \frac{\delta(t_k)^{2}}{\sigma^{2}_{\hat{B}}(t_k) + \sigma^{2}_{\mathrm{inst}}},
\end{equation}
donde $\boldsymbol{\theta}_{\mathrm{libre}} = \{\tau_h, \beta_0, \rho_{\mathrm{conv}}, T_A, T_{\mathrm{opt}}, T_H\}$ es el vector de parámetros libres del DEB sujetos a calibración (véase \cref{sec:cap4_calibracion}), $\sigma^{2}_{\hat{B}}(t_k)$ es la varianza predictiva del GP en el centro de la ventana, y $\sigma^{2}_{\mathrm{inst}}$ es la varianza instrumental del sensor NDIR (estimada en \num{625} ppm$^2$ a partir de la calibración de Fase I, correspondiente a un error estándar de \num{25} ppm). La ponderación inversa de la varianza tiene un efecto de \emph{down-weighting} automático: las ventanas donde el GP es incierto (alta $\sigma^{2}_{\hat{B}}$) contribuyen menos a la función de pérdida, protegiendo la calibración contra estimaciones de biomasa poco fiables \parencite{biagiDevelopmentMachineLearningbased2024}.

\subsection{Tratamiento del \ce{CH4} como residual de diagnóstico}
\label{subsec:cap4_tratamiento_ch4}

El metano no forma parte de la función de pérdida definida en la \cref{eq:cap4_funcion_perdida}, dado que el modelo DEB asume aerobiosis estricta y, por construcción, predice $\dot{m}_{\mathrm{CH_4}}^{\mathrm{DEB}} \equiv 0$. Esta exclusión no es una omisión, sino una decisión epistémica deliberada: el \ce{CH4} se reserva como variable de diagnóstico independiente, capaz de detectar desviaciones del régimen aeróbico que el modelo no puede ni debe reproducir.

El tratamiento del \ce{CH4} opera en dos niveles:

\paragraph{Nivel 1: Clasificador binario de anaerobiosis.}
La ecuación \eqref{eq:cap4_clasificador_ch4} del \cref{sec:cap4_arquitectura} define el indicador $\mathcal{A}(t) \in \{0, 1\}$ con umbral de \num{100} ppm. Este umbral se seleccionó por encima del nivel de ruido de fondo del sensor electroquímico de \ce{CH4} (\num{20}--\num{40} ppm de deriva) y por debajo de concentraciones biológicamente significativas asociadas a microambientes anaeróbicos en sustratos saturados \parencite{bekkerImpactSubstrateMoisture2021}. Cuando $\mathcal{A}(t) = 1$, el sistema emite una alerta operativa de nivel rojo, sugiriendo intervención inmediata (volteo del sustrato, ajuste de aireación).

\paragraph{Nivel 2: Residuo de coherencia modelo-sensor.}
Más allá de la detección binaria, el \ce{CH4} aporta información cualitativa sobre la validez del residuo $\delta(t_k)$. Si en una ventana de \num{30} min se observa simultáneamente:
\begin{itemize}
    \item $\mathcal{A}(t_k) = 1$ (presencia de \ce{CH4}), y
    \item $\delta(t_k) < -\num{50}$ ppm/min (el sensor de \ce{CO2} mide una pendiente sustancialmente menor que la predicha por el DEB),
\end{itemize}
el sistema interpreta que una fracción del flujo de carbono ha sido redirigida hacia la metanogénesis anaeróbica, en lugar de la respiración aeróbica modelada. En este escenario, el residuo $\delta(t_k)$ no se atribuye a error paramétrico del DEB, sino a una violación del supuesto de aerobiosis. Consecuentemente, dicha ventana se excluye de la optimización de $\boldsymbol{\theta}_{\mathrm{libre}}$ en la \cref{eq:cap4_funcion_perdida}, evitando que la calibración se corrompa por datos operativamente anómalos \parencite{rossiEstimatingDynamicsGreenhouse2024}.

\paragraph{Integración en el dashboard.}
La salida del módulo de pérdida se visualiza en el interfaz de decisión como una serie temporal de residuos $\delta(t_k)$, superpuesta con los eventos $\mathcal{A}(t_k) = 1$ marcados en color ámbar. Un residuo persistentemente negativo ($\delta < 0$ durante más de \num{3} ventanas consecutivas) sin detección de \ce{CH4} sugiere una subestimación sistemática de la biomasa por parte del GP o una deriva del sensor NDIR, activando una revisión de calibración. Por el contrario, un residuo positivo sostenido ($\delta > 0$) indica que el modelo DEB subestima la respiración real, lo que puede deberse a una temperatura efectiva del microambiente larval superior a la medida ambiental, o a una tasa de asimilación $a(T)$ subestimada por el kernel térmico \parencite{eriksenDynamicModellingFeed2022}.

La confrontación modelo-sensor, formalizada mediante la función de pérdida \eqref{eq:cap4_funcion_perdida} y el canal diagnóstico de \ce{CH4}, constituye el mecanismo de cierre del lazo híbrido. En la siguiente sección se establecen las métricas predictivas y los criterios estadísticos que permiten validar cuantitativamente el desempeño del sistema integrado.
